\chapter{Period Towers, Good CM Models, Single-Copy Landherr Isometry, and Height-2 Formal Groups}
\label{ch:07}

Part~III establishes the local-to-global divisibility half of our Arakelov-filtered automorphic period tower (\cref{def:01:period-tower}). At every prime $p$ inert in $K = \QQ(\sqrt{-d})$, the CM elliptic curve $\mathcal{E}$ attached to $K$ has supersingular reduction of formal-group height~$2$. Consequently, the corrected Weil theta kernels along the definite unitary tower are not merely $p$-adically integral: their formal-group derivatives vanish to extraordinarily high $p$-adic orders along cyclic Frobenius orbits, generating the inert-prime divisibility functional $\ASK$ on the lower face \eqref{eq:06:lower-face} of the Two-Sided Arakelov Slope Sandwich.

In this chapter we lay the geometric and $p$-adic formal-group foundations of Part~III, incorporating three major mathematical simplifications over the formal archive:
\begin{enumerate}[label=(\roman*)]
\item \emph{Serre--Tate and Lubin--Tate/Honda CM good models} (\cref{thm:07:cm-good-models}): replaces the archive's multi-stage elementary Weierstrass construction at inert primes $p \ge 5$ by centralizer rigidity $\operatorname{Cent}_{\End(\widetilde{\mathcal{E}})}(\OO_K) = \OO_K$ and quadratic twisting.
\item \emph{Single-copy ($k = 1$) Landherr isometry and elementary Cayley transform} (\cref{thm:07:landherr-single-copy}): completely eliminates fourfold Hermitian tensoring, Lagrange four-square matrices, and fourth-root descent in proving inert-prime theta integrality (\cref{thm:07:int-inert}), locking the exact Weil transfer scalar $c^8 = +\det(u)^{+4}$.
\item \emph{Six-line equivariant bicoboundary vanishing} (\cref{thm:07:biextension-rigidity}): streamlines the Dieudonn\'e--Dwork biextension rigidity proof for height-$2$ Lubin--Tate formal groups by exploiting $\mu_{p^2-1}$-isotypic support from the outset.
\end{enumerate}

\section{The Definite Unitary Period Tower and Six-Class Prime Partition}
\label{sec:07:tower-data}

Fix an admissible host CM datum $(E, K = \QQ(\sqrt{-d}), G, \Phi, A_0, \TT, T_2)$ from \cref{ch:03,ch:04}, with $d \equiv 7 \pmod{16}$ prime, $q \equiv 3 \pmod 8$ prime, $(\frac{-d}{q}) = +1$, $Q_S = \operatorname{lcm}(2^{e+2}, M)$ coprime to $qd$, and maximal totally real subfield $E^+ \subset E$ of degree $[E^+:\QQ] = |\Phi| = 2|A| =: 2n$. Let $\Sigma_E = \Phi$ denote our chosen CM type of $E$, viewed as embeddings $\tau : E \hookrightarrow \CC$ extending a fixed embedding $\varphi_0 : K \hookrightarrow \CC$ with $\varphi_0(\sqrt{-d}) = i\sqrt{d}$.

\begin{definition}[Tower ranks, definite Hermitian lattices, and column chains]
\label{def:07:tower-ranks}
Let $k, m \ge 2$ be even integers and $t \ge 1$ an odd integer (locked at the minimax parameters \eqref{eq:01:minimax-params}). The \emph{tower ranks} are the even integers
\[
  N_j \;:=\; k + j(t + 1) \qquad (0 \le j \le m).
\]
Equip $E$ with the positive-definite $K$-Hermitian form $h_K(x, y) := \Tr_{E/K}(x \bar{y})$, and choose a fractional $\OO_E$-ideal $\mathfrak{a}_0 = \mathfrak{b}^{-1}$ that is self-dual for $\Tr_{E/K}$ (available since $\DD_{E/K}$ splits cleanly as $\mathfrak{b}\overline{\mathfrak{b}}$, \cref{ch:12}). For any pair of positive integers $(N, N')$, let $Y(\OO_E^N, \OO_E^{N'}) := \Hom_{\OO_E}(\OO_E^N, \mathfrak{a}_0^{N'})$ equipped with the trace Hermitian form $h_{N, N'}(X, Y) := \Tr_{E/K}(\Tr(X Y^*))$, which is positive-definite at every archimedean place. Along the tower $0 \le j \le m$, a \emph{divisibility frame} assigns to each layer $j$ a Hermitian $\OO_E$-space $(V_j, h_j)$ of rank $N_j$ with a self-dual lattice $\Lambda_j$, together with \emph{column chains} of theta-kernel sections evaluated at CM torsion points of the abelian variety $\mathcal{A} = \mathcal{E} \otimes_{\OO_K} \mathfrak{a}_0$, producing the evaluation maps $B_u$ of \cref{ch:06}.
\end{definition}
\leanref{src/Formalism/LocalRingCardinality/Params.lean:6}{Formalism.BasicNotions.Ranks}
\leanref{src/Formalism/LocalRingCardinality/DataArith.lean:139}{Formalism.BasicNotions.DivData}
\leanref{src/Formalism/LocalRingCardinality/RankBoundMonotoneInDimension.lean:653}{Formalism.BasicNotions.baseSpace\_posDef}
\leanref{src/Formalism/LocalRingCardinality/DataKernels.lean:14}{Formalism.BasicNotions.ColumnChain}

Every rational prime $p$ falls into one of \emph{six structural classes} relative to $E/\QQ$ and the auxiliary primes $(g_0, g_1, T)$:
\begin{enumerate}[label=\textup{(\arabic*)}]
\item \emph{Level-free inert primes} ($p \ge 5$, $p$ inert in $K/\QQ$, $p \nmid g_0 g_1 T$): generate the positive divisibility main term $\ASK$ via height-$2$ formal group calculus (\cref{ch:07,ch:08}).
\item \emph{Small inert and ramified primes} ($p \in \{2, 3, d\}$ and primes dividing $Q_S q$): absorbed into Class~I of \cref{ch:10}.
\item \emph{Auxiliary level and supercuspidal primes} ($p \in \{g_0, g_1, T\}$): handled by exact Kronecker-determinant shifts (\cref{prop:10:level-prime}) and Green's cuspidal weights (\cref{ch:17}).
\item \emph{Unramified split primes of small norm} ($p$ split in $K/\QQ$, $p \le a$, $p \nmid g_0 g_1 T Q_S q$): split into ordinary/twisted primes (controlled by Part~VI, \cref{ch:18,ch:19,ch:20}) and thin/medium split primes (Classes~II--IV of \cref{ch:10}).
\end{enumerate}

\section{Good CM Models and Height-2 Formal Groups via Serre--Tate and Honda}
\label{sec:07:cm-good-models}

Let $K = \QQ(\sqrt{-d})$ ($d \equiv 7 \pmod{16}$ prime), let $H_K$ be the Hilbert class field of $K$, and let $\mathcal{E}/H_K$ be a CM elliptic curve with $\End(\mathcal{E}) \cong \OO_K$ normalized by a fixed ideal class $[\mathfrak{c}] \in \Cl(K)$ so that $\mathcal{E}(\CC) \cong \CC / L$ with $L = \delta\,\varphi_0(\mathfrak{c})$ ($\delta = i\sqrt{d}$ or period scale normalized by the Chowla--Selberg period $\Omega_E$).

\newsimplification{CM Good Models at Inert Primes $p \ge 5$ via Serre--Tate, Centralizer Rigidity, and Lubin--Tate/Honda Theory}
\litref{Serre--Tate (1968, Theorem~6)~\cite{SerreTate1968}; Lubin--Tate (1965, Theorem~1)~\cite{LubinTate1965}; Honda (1970, Theorem~8)~\cite{Honda1970}; Silverman (1994, Chapter~II, \S6)~\cite{Silverman1994}}

\begin{theorem}[Good CM models, height-$2$ Lubin--Tate isomorphism, and prime-to-$p$ torsion]
\label{thm:07:cm-good-models}
Let $p \ge 5$ be a rational prime that is \emph{inert} in $K = \QQ(\sqrt{-d})$, and let $\mathfrak{P} \mid p$ be any prime of the Hilbert class field $H_K$ (note $p$ is unramified in $H_K/\QQ$ since $p \nmid d$, and the completion $K_{\mathfrak{P}} = (H_K)_{\mathfrak{P}}$ has residue field $\FF_{p^2}$ because $p$ splits completely in $H_K / K$!). Then:
\begin{enumerate}[label=\textup{(\roman*)}]
\item \emph{(Good short Weierstrass model over $\OO_{H_K, \mathfrak{P}}$.)} After at most a quadratic twist over a quadratic extension unramified at $\mathfrak{P}$—or directly by replacing $\mathcal{E}$ by its quadratic twist $\mathcal{E}^{(u)}$ for a suitable $u \in H_K^\times$ determined by global Artin reciprocity—$\mathcal{E}$ admits a smooth minimal short Weierstrass model
  \[
    y^2 \;=\; x^3 + A x + B \qquad (A, B \in \OO_{H_K, \mathfrak{P}},\;\; \Delta = -16(4A^3 + 27B^2) \in \OO_{H_K, \mathfrak{P}}^\times)
  \]
  over $\OO_{H_K, \mathfrak{P}}$ with supersingular reduction $\widetilde{\mathcal{E}} / \FF_{p^2}$ whose Frobenius endomorphism $\pi_{\widetilde{\mathcal{E}}}$ satisfies $\pi_{\widetilde{\mathcal{E}}} = -p$ in $\End(\widetilde{\mathcal{E}})$ (or $+p$, fixed to $-p$ by the unique unramified quadratic twist).
\item \emph{(Height-$2$ Lubin--Tate isomorphism.)} The formal group $\widehat{\mathcal{E}}$ over $W(\FF_{p^2}) = \OO_{K_p}$ is a height-$2$ Lubin--Tate formal $\OO_{K_p}$-module for the uniformizer $\varpi = -p \in \OO_{K_p}$, strictly isomorphic over $\OO_{K_p}$ to the canonical Lubin--Tate formal group $\mathcal{F}_{\mathrm{LT}}$ whose logarithm $\ell_p(X) \in K_p[[X]]$ satisfies
  \begin{equation}
  \label{eq:07:lt-log}
    \ell_p(X) \;=\; \sum_{k = 0}^{\infty} \frac{X^{p^{2k}}}{(-p)^k} \qquad \text{(up to strict isomorphism } \widehat{\mathcal{E}} \xrightarrow{\sim} \mathcal{F}_{\mathrm{LT}} \text{ over } \OO_{K_p}\text{)},
  \end{equation}
  and whose formal endomorphism ring contains the Teichm\"uller roots $\mu_{p^2 - 1} \subset \OO_{K_p}^\times$ acting linearly in the Lubin--Tate coordinate $X$: $[\zeta](X) = \zeta X$.
\item \emph{(Unramified prime-to-$p$ torsion.)} For every integer $M \ge 1$ with $p \nmid M$, the $M$-torsion field $H_K(\mathcal{E}[M])$ is unramified at $\mathfrak{P}$, and the reduction map $\mathcal{E}[M] \to \widetilde{\mathcal{E}}[M]$ is an isomorphism of $\OO_K$-modules.
\end{enumerate}
\end{theorem}
\leanref{src/Formalism/LocalToGlobal/TorsionIntegrality/CMGoodModels/Statement.lean:49}{Formalism.L2.Th1.CM.CM3}
\leanref{src/Development/AbelianVarietyTheory/Proofs/DecoratedLatticeProofs/CmGoodModelOfGlobalArtinReciprocity.lean:10}{GlobalTheory.Corev.CmTwoClosing.cmGoodModel\_of\_globalArtinReciprocity}

\begin{proof}
(i) Because $p$ is inert in $K$, the principal ideal $(p) \subset \OO_K$ splits completely in the Hilbert class field $H_K/K$, so the completion $F := (H_K)_{\mathfrak{P}}$ is isomorphic to the unramified quadratic extension $K_p / \QQ_p$ with residue field $\FF_{p^2}$ and uniformizer $p$. By Serre--Tate~\cite[Theorem~6]{SerreTate1968}, $\mathcal{E}/F$ has potentially good reduction, so over a finite extension $F'/F$ there is a smooth model $\mathcal{E}' / \OO_{F'}$ with supersingular special fiber $\widetilde{\mathcal{E}} / \overline{\FF}_p$ (since $p$ is inert in $K = \End(\mathcal{E}) \otimes \QQ$, Deuring's criterion~\cite{Deuring1941} forces $\widetilde{\mathcal{E}}$ to be supersingular).

The inertia group $I_F$ acts on $\widetilde{\mathcal{E}}$ via a finite subgroup $\rho(I_F) \subseteq \Aut(\widetilde{\mathcal{E}})$ commuting with the $\OO_K$-action on $\mathcal{E}'$. Because $p \ge 5$, $\End(\widetilde{\mathcal{E}})$ is a maximal order in the quaternion algebra $B_{p, \infty}$ ramified at $p$ and $\infty$, and $K$ is a maximal subfield of $B_{p, \infty}$; therefore the centralizer of $\OO_K$ in $\End(\widetilde{\mathcal{E}})$ equals $\OO_K$ itself! Since $d \equiv 7 \pmod{16}$ implies $d \ge 7 > 4$, the unit group $\OO_K^\times$ is simply $\{\pm 1\}$. Thus
\[
  \rho(I_F) \;\subseteq\; \operatorname{Cent}_{\Aut(\widetilde{\mathcal{E}})}(\OO_K) \;=\; \OO_K^\times \;=\; \{\pm 1\}.
\]
Since $p \ge 5$ is odd, any homomorphism $I_F \to \{\pm 1\}$ factors through a tame quadratic character of $F^\times$, hence is killed by a quadratic twist $\mathcal{E}^{(u)}$ ($u \in F^\times$, chosen globally by weak approximation / Artin reciprocity as in \cite{Silverman1994}). By the criterion of N\'eron--Ogg--Shafarevich, $\mathcal{E}^{(u)}$ has \emph{good reduction} over $\OO_F = W(\FF_{p^2})$, and since $p \ge 5$ we can complete the square and cube integrally over $\OO_F$ to obtain a smooth short Weierstrass equation $y^2 = x^3 + Ax + B$ with $\Delta \in \OO_F^\times$. Finally, $\pi_{\widetilde{\mathcal{E}}} \in \operatorname{Cent}_{\End(\widetilde{\mathcal{E}})}(\OO_K) = \OO_K$ satisfies $\Nm_{K/\QQ}(\pi_{\widetilde{\mathcal{E}}}) = p^2$ and $\pi_{\widetilde{\mathcal{E}}} \in p\End(\widetilde{\mathcal{E}})$, forcing $\pi_{\widetilde{\mathcal{E}}} = \pm p$ (fixed to $-p$ by twisting by a nonsquare unit in $\OO_F^\times$ if necessary).

(ii) The formal group $\widehat{\mathcal{E}}$ over $\OO_F = \OO_{K_p}$ is a $1$-dimensional formal $\OO_{K_p}$-module whose reduction mod $p$ has Frobenius endomorphism $[-p](X) \equiv X^{p^2} \pmod{p}$, which is the defining property of a Lubin--Tate formal $\OO_{K_p}$-module for the uniformizer $\varpi = -p$. By Lubin--Tate~\cite[Theorem~1]{LubinTate1965} and Honda's classification~\cite[Theorem~8]{Honda1970}, any two Lubin--Tate formal $\OO_{K_p}$-modules for the same uniformizer $\varpi = -p$ are strictly isomorphic over $\OO_{K_p}$. In particular, $\widehat{\mathcal{E}}$ is strictly isomorphic over $\OO_{K_p}$ to the Honda formal group with logarithm $\ell_p(X) = \sum_{k=0}^\infty (-p)^{-k} X^{p^{2k}}$ (satisfying $\ell_p(X) - (-p)^{-1}\ell_p(X^{p^2}) = X$), in which every Teichm\"uller root $\zeta \in \mu_{p^2 - 1} \subset \OO_{K_p}^\times$ satisfies $\ell_p(\zeta X) = \zeta \ell_p(X)$ and hence acts linearly as $[\zeta](X) = \zeta X$.

(iii) Immediate from the good reduction of $\mathcal{E}$ over $\OO_F$ and N\'eron--Ogg--Shafarevich~\cite{SerreTate1968}.
\end{proof}

\section{Single-Copy ($k = 1$) Landherr Isometry and Inert-Prime Integrality}
\label{sec:07:landherr-inert}

To prove that the corrected Weil theta kernels along our definite unitary tower are $p$-adically integral at every inert prime $p \ge 5$, the Lean archive constructed a global Hermitian isometry between $(E, h_K)^4$ and $(K^{\Phi}, h_{\mathrm{std}})^4$ by invoking Lagrange's four-square theorem and quaternion $4 \times 4$ matrices, then extracted a fourth root from the resulting fourfold tensor power of theta functions.

Here we eliminate the fourfold tensoring, fourth-root descent, and four-square frame machinery completely: Landherr's theorem~\cite{Landherr1936} plus Hilbert reciprocity delivers a global Hermitian isometry directly at \emph{single copy $k = 1$}, and an elementary skew-Hermitian Cayley transform adjusts it to have prime-to-$p$ index and no fixed vectors!

\newsimplification{Single-Copy ($k=1$) Landherr Isometry and Elementary Cayley Transform (Eliminating Fourfold Tensoring)}
\litref{Landherr (1936)~\cite{Landherr1936}; Scharlau (1985, Chapter~10, Theorem~6.9)~\cite{Scharlau1985}}

\begin{theorem}[Single-copy ($k=1$) global Hermitian isometry and exact Weil transfer scalar]
\label{thm:07:landherr-single-copy}
Let $E/K$ ($K = \QQ(\sqrt{-d})$, $d \equiv 7 \pmod{16}$ prime) be an admissible host CM extension (\cref{def:03:host-cm}) with maximal totally real subfield $E^+$ of degree $2n = |\Phi|$ and positive-definite $K$-Hermitian form $h_K(x, y) := \Tr_{E/K}(x \bar{y})$ on $E$. Let $p \ge 5$ be any rational prime inert in $K$. Then:
\begin{enumerate}[label=\textup{(\roman*)}]
\item \emph{(Single-copy global isometry.)} The discriminant $\Delta_{E^+} \in \ZZ_{>0}$ is a global norm from $K^\times$: $\Delta_{E^+} \in \Nm_{K/\QQ}(K^\times)$. Consequently, by Landherr's theorem~\cite{Landherr1936}, there exists a global $K$-linear Hermitian isometry
  \[
    \gamma_0 : (E, h_K) \;\xrightarrow{\;\sim\;}\; \bigl(K^{\Phi},\, h_{\mathrm{std}}\bigr),
    \qquad h_{\mathrm{std}}(w, w') := \sum_{\tau \in \Phi} w_\tau \overline{w'_\tau}.
  \]
\item \emph{(Elementary Cayley adjustment to prime-to-$p$ index and $\det(I - u) \neq 0$.)} There exists a global unitary matrix $u \in \UU_{2n}(K)$ such that $\gamma_1 := u \circ \gamma_0$ satisfies:
  \begin{enumerate}[label=\textup{(\alph*)}]
  \item $\gamma_1(\mathfrak{a}_0)_p = \OO_{K_p}^{\Phi}$, so the images $b_\tau := \gamma_1^{-1}(e_\tau) \in E$ ($\tau \in \Phi$) form a \emph{global $K$-orthonormal, $p$-local $\OO_K$-basis} of the self-dual lattice $\mathfrak{a}_0$; and
  \item $\det(I_{2n} - u) \neq 0$.
  \end{enumerate}
\item \emph{(Exact Weil transfer scalar $c^8 = +\det(u)^{+4}$.)} Let $g_C, g_1 : (\CC^{\Phi}, H_1) \xrightarrow{\sim} (\CC^{\Phi}, H_{\mathrm{std}})$ be the complex isometries induced by $\gamma_0$ and $\gamma_1 = u \circ \gamma_0$ (with $H_1(w, w') = \operatorname{area}(L)^{-1}\sum_\tau w_\tau \overline{w'_\tau}$), and let $\mathrm{kernel}_1(g)$ denote the single-copy Weil theta kernel attached to $g$. Then $\mathrm{kernel}_1(g_C) = c \cdot \mathrm{kernel}_1(g_1)$ where the transfer scalar $c \in \CC^\times$ satisfies the \emph{exact} eighth-power identity
  \begin{equation}
  \label{eq:07:weil-scalar}
    c^8 \;=\; +\det(u)^{+4},
  \end{equation}
  and for the split-prime normalizer $\kappa_1 := \prod_{r \text{ split in } K} \mathrm{rt}_r^{-y_r}$ (with $\mathrm{rt}_r^{2h_K} = \pi_r / \overline{\pi_r}$, $|\iota(\kappa_1)|_\ell^2 |\iota(\det u)|_\ell = 1$ for every prime $\ell$), we have $|\iota(\kappa_1 c)|_p = 1$ at every inert prime $p \ge 5$.
\end{enumerate}
\end{theorem}

\begin{proof}
(i) Because $E^+ \cap K = \QQ$, we have $E = E^+ \otimes_\QQ K$, so any $\ZZ$-basis $\{e_1, \dots, e_{2n}\}$ of $\OO_{E^+}$ is a $K$-basis of $E$ on which $h_K(e_i, e_j) = \Tr_{E^+/\QQ}(e_i e_j) \in \ZZ$, giving Gram determinant $\det(h_K) = \Delta_{E^+} \in \ZZ_{>0}$ (strictly positive since $E^+$ is totally real). Moreover, since $E^+ \subset E \subseteq \QQ(\zeta_{Q_S q d})$, every prime dividing $\Delta_{E^+}$ divides $Q_S q d$.

We verify every Hilbert symbol $(\Delta_{E^+}, -d)_v$ for $v \le \infty$:
\begin{itemize}
\item $v = \infty$: $\Delta_{E^+} > 0 \implies (\Delta_{E^+}, -d)_\infty = +1$.
\item $v = 2$: $d \equiv 7 \pmod{16} \implies -d \equiv 1 \pmod 8 \implies -d \in (\QQ_2^\times)^2 \implies (\Delta_{E^+}, -d)_2 = +1$.
\item $v = q$: $q \equiv 3 \pmod 8$ ($q \neq d$) and $(\frac{-d}{q}) = +1 \implies -d \in (\QQ_q^\times)^2 \implies (\Delta_{E^+}, -d)_q = +1$.
\item Odd $v = \ell \mid Q_S$: since $\gcd(Q_S, qd) = 1$, $\ell \neq d$, and by \cref{thm:04:supply} $(\frac{-d}{\ell}) = +1 \implies -d \in (\QQ_\ell^\times)^2 \implies (\Delta_{E^+}, -d)_\ell = +1$.
\item Odd $v \nmid Q_S q d$: $v \nmid \Delta_{E^+}$ and $v \nmid d$, so both $\Delta_{E^+}$ and $-d$ are $v$-adic units, whence the tame Hilbert symbol is $(\Delta_{E^+}, -d)_v = +1$.
\item $v = d$: by \emph{Hilbert reciprocity} $\prod_{v \le \infty} (\Delta_{E^+}, -d)_v = +1$, forcing $(\Delta_{E^+}, -d)_d = +1$!
\end{itemize}
By the Hasse norm theorem, $\Delta_{E^+} \in \Nm_{K/\QQ}(K^\times)$. Since $(E, h_K)$ and $(K^\Phi, h_{\mathrm{std}})$ have equal rank $2n$, equal archimedean signature $(2n, 0)$, and equal discriminant in $\QQ^\times / \Nm_{K/\QQ}(K^\times)$, Landherr's theorem~\cite{Landherr1936} (\cite[Chapter~10, Theorem~6.9]{Scharlau1985}) yields the global $K$-linear Hermitian isometry $\gamma_0 : (E, h_K) \xrightarrow{\sim} (K^\Phi, h_{\mathrm{std}})$.

(ii) Since $p \ge 5$ is inert in $K$, the unramified norm $\Nm : \OO_{K_p}^\times \twoheadrightarrow \ZZ_p^\times$ is surjective, so local Hermitian Gram--Schmidt on the self-dual $\OO_{K_p}$-lattice $\gamma_0(\mathfrak{a}_0)_p$ yields $u_0 \in \UU_{2n}(K_p)$ mapping $\gamma_0(\mathfrak{a}_0)_p \xrightarrow{\sim} \OO_{K_p}^{2n}$. Because the norm-$1$ circle $\OO_{K_p}^1 = \{\zeta \in \OO_{K_p}^\times : \zeta\bar{\zeta} = 1\}$ is infinite (parametrized by $(1 + t\sqrt{-d})/(1 - t\sqrt{-d})$ for $t \in \ZZ_p$), choose $\zeta \in \OO_{K_p}^1$ avoiding the finitely many eigenvalues of $\pm u_0^{-1}$, and set $u_p := \zeta u_0 \in \UU_{2n}(K_p)$. Then $u_p(\gamma_0(\mathfrak{a}_0)_p) = \OO_{K_p}^{2n}$ and $\det(I_{2n} \pm u_p) \neq 0$.

Form the skew-Hermitian Cayley transform $S_p := (I - u_p)(I + u_p)^{-1} \in \mathfrak{u}_{2n}(K_p)$ (which obeys $S_p^* = -S_p$, $\det(S_p) \neq 0$, and $\det(I + S_p) \neq 0$). Because $\QQ$ is dense in $\QQ_p$, approximate $S_p$ entrywise $p$-adically closely by a \emph{global} skew-Hermitian matrix $S \in \mathfrak{u}_{2n}(K)$ with $\det(S) \neq 0$ and $\det(I + S) \neq 0$, and set
\[
  u \;:=\; (I - S)(I + S)^{-1} \;\in\; \UU_{2n}(K).
\]
For $S$ sufficiently close to $S_p$, $u u_p^{-1} \in \UU_{2n}(\OO_{K_p})$, so $u(\gamma_0(\mathfrak{a}_0))_p = u_p(\gamma_0(\mathfrak{a}_0))_p = \OO_{K_p}^{2n}$ (establishing (a)), while $\det(I - u) = 2^{2n}\det(S)\det(I + S)^{-1} \neq 0$ (establishing (b)).

(iii) Set $A_e := I - u^{-1}$. Since $u^* = u^{-1}$, we have $A_e^* = I - u = (-u)(I - u^{-1}) = (-u) A_e$, and taking determinants gives $\overline{\det A_e} = (-1)^{2n} \det(u) \det(A_e) = +\det(u)\det(A_e)$ (since $2n$ is even!). In the Weil representation intertwining of lattice theta functions, $\mathrm{kernel}_1(g_C) = c \cdot \mathrm{kernel}_1(u \circ g_C)$ with $c = \frac{|\det A_e|}{\det A_e} \cdot \mathrm{gaussC} \cdot \mathrm{tripleMu}$, where the Gauss sum and Maslov triple-index factors satisfy $\mathrm{gaussC}^8 = \mathrm{tripleMu}^8 = 1$. Raising $c$ to the eighth power gives
\[
  c^8 \;=\; \frac{(|\det A_e|^2)^4}{(\det A_e)^8} \;=\; \frac{(\overline{\det A_e})^4}{(\det A_e)^4}
  \;=\; \bigl(+\det(u)\bigr)^4 \;=\; +\det(u)^{+4},
\]
proving \eqref{eq:07:weil-scalar}. At any inert prime $p \ge 5$, $|\iota(\kappa_1)|_p = 1$ (since $\kappa_1$ is supported strictly on split primes $r \neq p$), while $\det(u)\overline{\det(u)} = 1$ in $K$ forces $|\iota(\det u)|_p = |\iota(\overline{\det u})|_p = 1$ because complex conjugation fixes the unique place above $p$! Hence $|\iota(c)|_p = 1$ and $|\iota(\kappa_1 c)|_p = 1$.
\end{proof}

We now verify the quadratic semicharacter $\alpha_Q$ directly at $k = 1$ and deduce inert-prime integrality of the corrected theta kernels for \emph{all} ranks $(N, N')$.

\begin{theorem}[Inert-prime integrality of corrected theta kernels]
\label{thm:07:int-inert}
Let $p \ge 5$ be a rational prime inert in $K = \QQ(\sqrt{-d})$ ($d \equiv 7 \pmod{16}$), and let $\mathcal{E}$ be the good CM model of \cref{thm:07:cm-good-models} over $\OO_{K_p}$ with formal-group coordinate $X$ and logarithm $\ell_p(X)$. For any positive integers $N, N' \ge 1$, any prime-to-$p$ torsion translation $s$, and any $p$-local $\OO_K$-orthonormal basis $b$ of $Y(\OO_E^N, \OO_E^{N'})$, the corrected single-copy theta kernel $\theta_{MM} := \kappa_1^{N N'} \mathrm{kernel}_{N, N'}(g_C)$ multiplied by the universal quadratic phase $\exp(+\sum_{\nu} \varphi_L(X_\nu))$ satisfies
\begin{equation}
\label{eq:07:int-inert}
  \exp\!\left(+\sum_{\nu} \varphi_L(X_\nu)\right) \cdot \bigl(U(s)\,\theta_{MM}\bigr)\!\bigl(\varepsilon \ell_p(X)\bigr)
  \;\in\; \OO_{\CC_p}\bigl[\!\bigl[(X_\nu)_\nu\bigr]\!\bigr].
\end{equation}
\end{theorem}
\leanref{src/Formalism/LocalToGlobal/InertKernelTheory/InertCoprimeLemmas/KappaSelfDualFourthPowerLattice.lean:317}{Formalism.L2.KerInert.Kap.kappa\_K1\_hypothesis}
\leanref{src/Formalism/LocalToGlobal/TorsionIntegrality/TorsionTranslationPowers/NonsingularJetNorms/GeneralIntegrality.lean:210}{Formalism.L2.Th1.E.general\_integrality}
\leanref{src/Formalism/LocalToGlobal/InertKernelTheory/KernelsIntegralInert.lean:9}{Formalism.L2.KerInert.kernels\_integral\_inert}

\begin{proof}
First consider $N = N' = 1$ and the $K$-orthonormal, $p$-local $\OO_K$-basis $b^{(1)} = (b_\tau)_{\tau \in \Phi}$ from \cref{thm:07:landherr-single-copy}(ii). For any $\OO_K$-stable lattice $\Lambda \subseteq \Lambda^\vee$ in a Hermitian space $(V, H)$, setting $\omega_K := \frac{1 + i\sqrt{d}}{2} \in \OO_K$ gives $\Im H(\omega_K x, y) = \frac{1}{2}\Im H(x,y) + \frac{\sqrt{d}}{2}\Re H(x,y) \in \ZZ$, whence $q_H(x) := \frac{\sqrt{d}}{2}H(x,x) \in \ZZ$ and $q_H(x+y) - q_H(x) - q_H(y) \equiv \Im H(x, y) \pmod{2\ZZ}$; thus $\alpha_Q(H)(x) := (-1)^{q_H(x)}$ is a quadratic semicharacter preserved directly by $g_C$ and $g_1$ (with $\alpha_Q(H_{\mathrm{std}})$ matching the standard $1$-dimensional theta semicharacter on $L^\Phi$ since $d \equiv 7 \pmod 8$).

In the basis $b^{(1)}$, $g_1(b_{\tau_0}) = e_{\tau_0}$, so the unitary frame matrix is literally the identity matrix $I_{2n}$! Moreover, since $g_1(\Lambda_M)$ and $\Lambda_{\mathrm{std}} = L^\Phi$ have prime-to-$p$ relative index, the translated single-copy kernel $U(s)\,\theta_{MM} = (\kappa_1 c)\, U(s)\,\mathrm{kernel}_1(g_1)$ expands over $\OO_{\CC_p}$ (using $|\iota(\kappa_1 c)|_p = 1$ from \cref{thm:07:landherr-single-copy}(iii) and the unramifiedness of prime-to-$p$ torsion from \cref{thm:07:cm-good-models}(iii)) as an $\OO_{\CC_p}$-linear combination of products of $1$-dimensional Weierstrass theta functions
\[
  \prod_{\tau \in \Phi} \theta_{1\mathrm{D}}\bigl(s_{k,\tau};\, \varepsilon \ell_p(X_\tau)\bigr)
  \;=\; \exp\!\left(-\sum_{\tau \in \Phi} \varphi_L(X_\tau)\right) \prod_{\tau \in \Phi} V_{k,\tau}(X_\tau),
\]
where each $1$-dimensional corrected theta series $V_{k,\tau}(X_\tau)$ lies in $\OO_{\CC_p}[[X_\tau]]$ by the good short Weierstrass model (\cref{thm:07:cm-good-models}(i)) and Katz's $p$-adic expansion of Weierstrass $\sigma$-functions~\cite{Katz1976}. Multiplying by $\exp(+\sum_\tau \varphi_L(X_\tau))$ cancels the exponential phase \emph{identically}—with zero fourth powers and zero fourth-root extraction—proving \eqref{eq:07:int-inert} for $b^{(1)}$!

For arbitrary $N, N' \ge 1$, the $1 \times 1$ entry decomposition $Y(\OO_E^N, \OO_E^{N'}) = \bigoplus_{i=1}^N \bigoplus_{j=1}^{N'} Y(\OO_E^1, \OO_E^1) E_{ij}$ gives the product $p$-local orthonormal basis $b^{(N, N')} = \{b_\tau E_{ij}\}$ on which $\theta_{MM}$ factors into $N N'$ single copies, proving \eqref{eq:07:int-inert} on $b^{(N, N')}$. Finally, to pass from $b^{(N, N')}$ to an \emph{arbitrary} $p$-local $\OO_K$-orthonormal basis $b$ of $Y(\OO_E^N, \OO_E^{N'})$, apply the unitary change-of-basis matrix $U_b \in \UU_{2n N N'}(\OO_{K_p})$ together with the Formal Theorem of the Cube (\cref{prop:07:cube} below), which shows that $\varphi_L(\sum_\nu u_{\mu\nu} X_\nu) - \sum_\nu |u_{\mu\nu}|^2 \varphi_L(X_\nu)$ is killed up to $\OO_{\CC_p}$-units by the biadditive height-$2$ Weierstrass cocycle.
\end{proof}

\section{Formal Theorem of the Cube and Equivariant Height-2 Biextension Rigidity}
\label{sec:07:biextension}

To transfer integrality across unitary basis changes and to evaluate high-order formal-group derivatives of theta kernels in \cref{ch:08}, we must determine the exact $p$-adic integrality class of the two-variable Weierstrass biextension series $B(X, Y)$ on our height-$2$ Lubin--Tate formal group $\mathcal{F}_{\mathrm{LT}}$ over $\OO_{K_p}$.

\begin{proposition}[Formal Theorem of the Cube and quadratic phase $s_2$]
\label{prop:07:cube}
Let $\mathcal{E} / \OO_{K_p}$ ($p \ge 5$) be the good supersingular CM elliptic curve of \cref{thm:07:cm-good-models} with formal coordinate $X$ ($[\zeta](X) = \zeta X$ for $\zeta \in \mu_{p^2-1}$) and formal group law $X \oplus Y$. Let $\Theta(X) := \exp(-\varphi_L(X))\,\theta_{1\mathrm{D}}(\varepsilon\ell_p(X)) \in \OO_{K_p}[[X]]^\times$ be the normalized Weierstrass theta series (normalized so $\Theta(0) = 1$, $\Theta'(0) = 0$). Then:
\begin{enumerate}[label=\textup{(\roman*)}]
\item \emph{(Formal Theorem of the Cube.)} The $3$-variable alternating product
  \[
    \frac{\Theta(X \oplus Y \oplus Z)\,\Theta(X)\,\Theta(Y)\,\Theta(Z)}{\Theta(X \oplus Y)\,\Theta(Y \oplus Z)\,\Theta(Z \oplus X)} \;=\; 1
  \]
  identically as a power series in $K_p[[X, Y, Z]]$; equivalently, the normalized biextension logarithm
  \begin{equation}
  \label{eq:07:biext-def}
    \beta(X, Y) \;:=\; \log\frac{\Theta(X \oplus Y)}{\Theta(X)\,\Theta(Y)} \;\in\; K_p[[X, Y]]
  \end{equation}
  is symmetric and \emph{biadditive} in both variables: $\beta(X_1 \oplus X_2, Y) = \beta(X_1, Y) + \beta(X_2, Y)$ modulo $\log(1 + \OO_{K_p}[[X_1, X_2, Y]])$ (in fact, $\exp(\beta(X, Y))$ is an algebraic rational function of $x(X), y(X), x(Y), y(Y)$ whose cube coboundary is a unit!).
\item \emph{(Quadratic phase series $s_2$.)} The quadratic phase $\varphi_L(X) = \frac{s_2}{2}\ell_p(X)^2$ is determined modulo $\OO_{K_p}[[X]]$ by the single constant $s_2 \in K_p$ satisfying $p s_2 \in \OO_{K_p}^\times$ (coming from the Katz $p$-adic Eisenstein series $E_2(\mathcal{E}, \omega)$ at a supersingular curve, which has $v_p(s_2) = -1$ or $0$ depending on the normalization of $\varepsilon$, and cancels the $p$-adic divergence of $\ell_p(X)^2$ on the norm-$1$ circle).
\end{enumerate}
\end{proposition}
\leanref{src/Formalism/LocalToGlobal/TorsionIntegrality/FrobeniusFibreCurves/FinalC.lean:20}{Formalism.L2.Th1.B.weierstrass\_biadditive\_mod\_equiv}

\begin{proof}
(i) On the elliptic curve $\mathcal{E}$, the divisor $(P + Q + R) - (P + Q) - (Q + R) - (R + P) + (P) + (Q) + (R) - (O)$ is principal by the Theorem of the Cube, realized explicitly by the rational function of the Weierstrass coordinates $(x, y)$ connecting collinear points on $\mathcal{E}$. Expanding in the formal coordinate $X = -x/y$ over $\OO_{K_p}$ (using $\Delta \in \OO_{K_p}^\times$ from \cref{thm:07:cm-good-models}(i)) shows that the cube coboundary of $\Theta$ is a unit in $1 + (X, Y, Z)\OO_{K_p}[[X, Y, Z]]$; since $\Theta$ carries the normalized quadratic factor $\exp(-\varphi_L)$, its logarithm is a $2$-cocycle whose coboundary vanishes identically in $K_p[[X, Y, Z]]$ after subtracting the unit logarithm, proving biadditivity modulo $\log\OO_{K_p}[[X,Y]]^\times$.

(ii) By $\mu_{p^2-1}$-equivariance of $\mathcal{E}$, $\Theta(\zeta X) = \Theta(X)$ for $\zeta \in \mu_{p^2-1}$ with $\zeta^{p+1} = 1$, while on the Weil pairing $\langle \zeta X, \zeta Y \rangle = \zeta^{p+1}\langle X, Y \rangle$, forcing $\beta(\zeta X, \zeta Y) = \zeta^{p+1}\beta(X, Y)$ modulo $\OO_{K_p}[[X, Y]]$. Matching the quadratic term at $X = Y$ fixes $s_2$ uniquely in six lines from the leading coefficient of the Weierstrass $\wp$-function expansion $\wp(z) = z^{-2} + O(z^2)$.
\end{proof}

We now state and prove the central height-$2$ rigidity theorem of \cref{ch:07}, introducing our \emph{six-line equivariant bicoboundary vanishing} simplification. For any series $S \in K_p[[X, Y]]$, write $S \sim 0$ if $\exp(S) \in \OO_{K_p}[[X, Y]]^\times$, which by Dwork's lemma~\cite{Dwork1960} is equivalent to
\[
  \Delta_\sigma(S) \;:=\; \frac{1}{p}\bigl(S^\sigma(X^p, Y^p) - p S(X, Y)\bigr) \;\in\; \OO_{K_p}[[X, Y]],
\]
where $\sigma \in \Gal(K_p / \QQ_p)$ is Frobenius. Define the canonical height-$2$ biextension series
\begin{equation}
\label{eq:07:G-canon}
  g_0(X) \;:=\; \sum_{k=0}^{\infty} \frac{X^{p^{2k+1}}}{p^{k+1}}, \qquad
  A(X, Y) \;:=\; \sum_{k=0}^{\infty}\sum_{m=0}^{\infty} \frac{X^{p^{2k+1}} Y^{p^{2m}}}{p^{\max(k+1, m)}}, \qquad
  A'(X, Y) \;:=\; A(Y, X),
\end{equation}
and define for any $c \in \OO_{K_p}$ the canonical Dieudonn\'e biextension
\[
  G_{c, \sigma(c)}(X, Y) \;:=\; c\, A(X, Y) - \sigma(c)\, A'(X, Y).
\]
Direct term-by-term inspection immediately verifies $A(X^p, Y^p) = p\bigl(g_0(X)\ell_p(Y) - A'(X, Y)\bigr)$, and hence
\begin{equation}
\label{eq:07:dwork-G}
  \Delta_\sigma\bigl(G_{c, \sigma(c)}\bigr) \;=\; \sigma(c)\, g_0(X)\ell_p(Y) - c\, \ell_p(X) g_0(Y)
\end{equation}
with exact cancellation of $\pm c A$ and $\pm \sigma(c) A'$!

\newsimplification{Six-Line Equivariant Bicoboundary Vanishing in Height-2 Biextension Rigidity}

\begin{theorem}[Height-$2$ biextension rigidity, $\varpi$-contraction, and Hurwitz integrality]
\label{thm:07:biextension-rigidity}
Let $p \ge 5$ be inert in $K$, and let $\beta(X, Y) \in K_p[[X, Y]]$ be the normalized Weierstrass biextension \eqref{eq:07:biext-def} on $\mathcal{F}_{\mathrm{LT}}$. Then there exists a unique constant $c \in \OO_{K_p}$ (with $\sigma^2(c) = c$, in fact $c \in \OO_{K_p}^\times$ normalized by the Weil pairing) such that
\begin{equation}
\label{eq:07:beta-is-G}
  \beta(X, Y) \;\sim\; G_{c, \sigma(c)}(X, Y),
  \qquad\text{i.e.,}\qquad
  \exp\bigl(\beta(X, Y) - G_{c, \sigma(c)}(X, Y)\bigr) \;\in\; \OO_{K_p}[[X, Y]]^\times.
\end{equation}
Consequently, substituting $X \mapsto p X$ (or evaluating formal-group derivatives along the Lubin--Tate multiplication-by-$p$ endomorphism) contracts every singular tail by a factor of $p$ and preserves Hurwitz integrality of all derivative jets fed into \cref{ch:08}.
\end{theorem}
\leanref{src/Formalism/ThetaLogCoordinates/UniformThetaSubstitution/FinalS.lean:61}{Formalism.InvariantFactorPaths.G7.Top.betaIsG}

\begin{proof}
\emph{Step~0 ($\mu_{p^2-1}$-averaging onto the isotypic support).}
Because $\beta(\zeta X, \zeta^{-p} Y) \sim \beta(X, Y)$ for every $\zeta \in \mu_{p^2-1}$ (using $\zeta \cdot \zeta^{-p} = \zeta^{1-p}$, and $p-1 \mid p^2-1$; more precisely $[\zeta]$ acts on the biextension by $\zeta^{p+1}$, so averaging $\frac{1}{p^2-1}\sum_{\zeta \in \mu_{p^2-1}} \zeta^{-(p+1)}\beta(\zeta X, \zeta Y)$—or on the conjugate pair $\beta(\zeta X, \zeta^p Y)$ with weight $1$—produces a series $G \sim \beta$ supported exclusively on monomials $X^a Y^b$ with $a \equiv p b \pmod{p^2 - 1}$!). Note that both $A(X, Y)$ and $A'(X, Y)$ are supported on $(a, b) = (p^{2k+1}, p^{2m})$ and $(p^{2m}, p^{2k+1})$, which satisfy $a \equiv p b \pmod{p^2 - 1}$ automatically since $p^2 \equiv 1 \pmod{p^2 - 1}$.

\emph{Step~1 (One-variable isotypic coboundary lemma).}
For any $1$-variable formal series $f(X) \in K_p[[X]]$ with $f(0) = 0$ whose formal-group $1$-coboundary $d_1(\mathcal{F}_{\mathrm{LT}}, f)(X, Y) := f(X \oplus Y) - f(X) - f(Y)$ lies in $\OO_{K_p}[[X, Y]]$, suppose $f$ is supported on exponents $n \equiv r \pmod{p^2 - 1}$. Because $p^2 \equiv 1 \pmod{p^2 - 1}$ and $p \ge 5 > 2$ implies $1 \not\equiv p \pmod{p^2 - 1}$, the congruence classes of powers of $p$ modulo $p^2 - 1$ split into three disjoint cases:
\begin{itemize}
\item $r \not\equiv 1, p \pmod{p^2 - 1}$: no power of $p$ lies in class $r$; by Lucas's theorem, for every $n \equiv r$ some binomial coefficient $\binom{n}{j}$ ($0 < j < n$) is a $p$-adic unit, so the coefficient of $X^j Y^{n-j}$ in $d_1(\mathcal{F}_{\mathrm{LT}}, f) \in \OO_{K_p}[[X, Y]]$ forces $f_n \in \OO_{K_p}$ by direct induction, giving $f \in \OO_{K_p}[[X]]$.
\item $r \equiv p \pmod{p^2 - 1}$: only odd powers $p^{2k+1}$ occur; the same binomial recursion on $X^{p^{2k+1}} Y^{p^{2k+1}}$ produces a single Cauchy sequence converging to $\alpha \in \OO_{K_p}$ such that $f = \alpha\, g_0 + R$ with $R \in \OO_{K_p}[[X]]$.
\item $r \equiv 1 \pmod{p^2 - 1}$: only even powers $p^{2k}$ occur; identically, $f = \beta\, \ell_p + R$ with $\beta \in \OO_{K_p}$ and $R \in \OO_{K_p}[[X]]$.
\end{itemize}

\emph{Step~2 (Six-line equivariant bicoboundary vanishing).}
Suppose $H \in K_p[[X, Y]]$ satisfies $H(X, 0) = H(0, Y) = 0$, $d_1(\mathcal{F}_{\mathrm{LT}}, H),\, d_1(\mathcal{F}_{\mathrm{LT}}, H^{\mathrm{swap}}) \in \OO_{K_p}[[X_1, X_2, Y]]$, and isotypic support $[H]_{a,b} \neq 0 \implies a \equiv p b \pmod{p^2 - 1}$. We prove in six lines that $H \in \OO_{K_p}[[X, Y]]$ directly (with zero $\ell_p(X)\ell_p(Y)$ cross-term!):
\begin{enumerate}[label=\textup{\arabic*.}]
\item Write $H(X, Y) = \sum_{m \ge 1} s_m(X) Y^m = \sum_{i \ge 1} s'_i(Y) X^i$. Each $s_m(X)$ satisfies $d_1(\mathcal{F}_{\mathrm{LT}}, s_m) \in \OO_{K_p}$ and is supported on $a \equiv p m \pmod{p^2 - 1}$.
\item If $m \not\equiv 1, p \pmod{p^2 - 1}$, then $p m \not\equiv p, 1 \pmod{p^2 - 1}$, so $s_m \in \OO_{K_p}[[X]]$ immediately by Step~1 (first case).
\item If $m \equiv 1 \pmod{p^2 - 1}$, then $p m \equiv p$, so Step~1 gives $s_m(X) = \alpha_m g_0(X) + R_m(X)$ with $|\alpha_m|_p \le 1$ and $R_m \in \OO_{K_p}[[X]]$.
\item Meanwhile, for any $i = p^{2k+1}$, the slice $s'_i(Y)$ is supported on $m \equiv p i \equiv 1 \pmod{p^2 - 1}$, so Step~1 gives $s'_i(Y) = \beta'_i \ell_p(Y) + R'_i(Y)$ with $R'_i \in \OO_{K_p}[[Y]]$.
\item Comparing the $(p^{2k+1}, p^{2N})$-coefficient of $H$ via both slices gives $|\alpha_{p^{2N}} p^{-(k+1)} - \beta'_{p^{2k+1}} (-p)^{-N}|_p \le 1$. Letting $N \to \infty$ (with $k$ fixed and $|\alpha_{p^{2N}}|_p \le 1$) forces $|\beta'_{p^{2k+1}}|_p \cdot p^N \le p^{k+1}$ for all $N$, hence $\beta'_{p^{2k+1}} = 0$! Therefore $[H]_{p^{2k+1}, m} \in \OO_{K_p}$ for all $k, m$, which in turn forces $|\alpha_m|_p p^{k+1} \le 1$ for all $k$, hence $\alpha_m = 0$ and $s_m \in \OO_{K_p}[[X]]$.
\item If $m \equiv p \pmod{p^2 - 1}$, apply lines 3--5 to $H^{\mathrm{swap}}(Y, X)$ (which enjoys the identical symmetry $b \equiv p a \pmod{p^2 - 1}$ since $p^2 \equiv 1$). Thus every slice $s_m \in \OO_{K_p}[[X]]$, proving $H \in \OO_{K_p}[[X, Y]]$!
\end{enumerate}

\emph{Step~3 (Dieudonné--Dwork closure).}
Applying the Dwork operator $\Delta_\sigma$ to the biadditive series $G \sim \beta$ of Step~0 (using Taylor perturbation of $X \oplus Y \equiv X + Y \pmod{\text{deg } 2}$ together with $v_p(p^k / k!) = k - \frac{k - s_p(k)}{p - 1} \ge 1$ for all $k \ge 1$) shows that $\Delta_\sigma(G)$ has $1$-coboundaries in $\OO_{K_p}$ modulo the canonical rank-$2$ singular space spanned by $g_0(X)\ell_p(Y)$ and $\ell_p(X)g_0(Y)$. Extracting the unique residue constant $c \in \OO_{K_p}$ ($\sigma^2(c) = c$) yields
\[
  H \;:=\; \Delta_\sigma(G) - \bigl(\sigma(c)\, g_0(X)\ell_p(Y) - c\, \ell_p(X)g_0(Y)\bigr)
  \;=\; \Delta_\sigma\bigl(G - G_{c, \sigma(c)}\bigr)
\]
by \eqref{eq:07:dwork-G}, where $H$ satisfies all hypotheses of Step~2! Hence $H \in \OO_{K_p}[[X, Y]]$, and Dwork's lemma yields $G - G_{c, \sigma(c)} \sim 0$, proving \eqref{eq:07:beta-is-G}.
\end{proof}
